%! The Nullspace of A: Solving Ax = 0 — Unit 3, Lesson 3.2 — Student Packet Cover Sheet
\documentclass[10pt]{article}
\usepackage{linalg-article}
\usepackage{linalg-boxes}
\usepackage{ltablex}
\keepXColumns

\begin{document}

% Full-bleed burgundy banner
\begin{tikzpicture}[remember picture, overlay]
  \fill[burgundy] (current page.north west)
    rectangle ([yshift=-0.9in]current page.north east);
\end{tikzpicture}
\vspace{-0.8in}

\begin{center}
  {\color{white}\textbf{\LARGE Linear Algebra}}\\[4pt]
  {\color{blushmid}\large\textbf{Unit 3 \quad The Four Fundamental Subspaces}}\\[6pt]
  {\color{white}\textbf{\large Lesson 3.2 \quad The Nullspace of $A$: Solving $A\mathbf{x}=\mathbf{0}$}}
\end{center}

\vspace{0.22in}
\namedateperiod
\vspace{0.18in}

\begin{learningtargetbox}
\small
\begin{enumerate}[leftmargin=1.5em, itemsep=5pt]
  \item \ldots{} describe the \textbf{nullspace} $N(A)$ --- all solutions of $A\mathbf{x}=\mathbf{0}$ --- and explain why it is a subspace of $\mathbb{R}^n$.
  \item \ldots{} solve $A\mathbf{x}=\mathbf{0}$ by \textbf{elimination}, telling \textbf{pivot columns} from \textbf{free columns}.
  \item \ldots{} build a \textbf{special solution} for each free variable and describe $N(A)$ as a point, line, or plane (there are $n-r$ free variables).
\end{enumerate}
\end{learningtargetbox}

\vspace{0.14in}

\begin{tocbox}
\small
\renewcommand{\arraystretch}{1.5}
\begin{tabularx}{\linewidth}{@{} c l X r @{}}
\rowcolor{blush}
\textbf{\#} & \textbf{Component} & \textbf{Description} & \textbf{Score} \\
\midrule
1 & Warm-Up        & Spiral review: verify $A\mathbf{x}=\mathbf{0}$, one elimination step, the subspace requirement & \blank{1.2cm} \\
2 & Guided Notes   & The nullspace, pivots vs.\ free columns, special solutions, and its geometry              & \blank{1.2cm} \\
3 & Group Activity & Find the Nullspace --- reducing, special solutions, and lines vs.\ planes                  & \blank{1.2cm} \\
4 & Exit Ticket    & Quick check on nullspaces, free variables, and what $N(A)=\{\mathbf{0}\}$ means           & \blank{1.2cm} \\
5 & Homework       & Nullspace practice, $C(A)$ vs.\ $N(A)$, and a complete-solution preview                    & \blank{1.2cm} \\
\midrule
  & \multicolumn{2}{r}{\textbf{Total}} & \blank{1.2cm} \\
\end{tabularx}
\end{tocbox}

\vspace{0.10in}

\begin{remindbox}
\small The \textbf{nullspace} $N(A)$ is every solution $\mathbf{x}$ of $A\mathbf{x}=\mathbf{0}$ --- the
``do-nothing'' inputs. It is a subspace of $\mathbb{R}^n$ (the columns' world). To find it, reduce $A$:
columns \emph{with} a pivot give \textbf{pivot variables}, columns \emph{without} one give \textbf{free
variables}. Set one free variable to $1$ (the rest to $0$) and solve back for a \textbf{special
solution}; $N(A)$ is all combinations of them. There are $n-r$ free variables, so $N(A)$ is a point
($n-r=0$), a line ($1$), or a plane ($2$) through the origin.
\end{remindbox}

\end{document}
